% ===================================================================
%  Generado por Tabula
%  Compílalo en Overleaf, o localmente con dos pasadas de pdflatex.
% ===================================================================
\documentclass[aspectratio=169,11pt]{beamer}
\usetheme{Boadilla}
\usefonttheme{serif}
% Boadilla es el tema de fábrica más limpio y, con serifas, es la puesta en
% escena de una defensa. En pantalla la letra es Termes; en el PDF, la serifa
% de fábrica de LaTeX. Para que coincidan, añade \usepackage{tgtermes}.
\definecolor{marinoPapel}{HTML}{F4F6FA}
\definecolor{marinoTinta}{HTML}{15243D}
\definecolor{marinoAzul}{HTML}{16345E}
\definecolor{marinoLaton}{HTML}{B08A32}
\setbeamercolor{background canvas}{bg=marinoPapel}
\setbeamercolor{normal text}{fg=marinoTinta}
\setbeamercolor{structure}{fg=marinoAzul}
\setbeamercolor{frametitle}{bg=marinoPapel,fg=marinoAzul}
\setbeamercolor{block title}{bg=white,fg=marinoAzul}
\setbeamercolor{block body}{bg=white}
\setbeamercolor{alerted text}{fg=marinoLaton!65!marinoTinta}
\setbeamerfont{frametitle}{series=\bfseries}
% El trozo corto de latón sobre la regla del título, el canto de color de las
% cajas y la barra de sección son de Tabula: Beamer no los trae y no salen
% sin reescribir las plantillas de frametitle y de bloque.
\usepackage[utf8]{inputenc}
\usepackage[T1]{fontenc}
% Español: se carga solo si babel-spanish está instalado (Overleaf lo trae).
\IfFileExists{spanish.ldf}{\usepackage[spanish,es-nodecimaldot]{babel}}{}
\usepackage{amsmath,amssymb}
\usepackage[version=4]{mhchem}
\definecolor{erlenfondo}{HTML}{F4F6FA}
\definecolor{erlenacento}{HTML}{B08A32}
\usepackage{tikz}
\usetikzlibrary{arrows.meta, calc}
\usepackage{pgfplots}
\pgfplotsset{compat=1.18}
\definecolor{serieA}{HTML}{0072B2}
\definecolor{serieB}{HTML}{E69F00}
\definecolor{serieC}{HTML}{009E73}
\definecolor{serieD}{HTML}{D55E00}
\definecolor{serieE}{HTML}{CC79A7}
\definecolor{serieF}{HTML}{56B4E9}
\renewcommand{\figurename}{Figura}
\renewcommand{\tablename}{Tabla}
\usepackage{siunitx}
\sisetup{per-mode=symbol, range-units=single}
\setbeamertemplate{navigation symbols}{}
% Algunos temas (Metropolis entre ellos) insertan solos una diapositiva al
% empezar cada sección. Tabula ya emite la suya, así que se desactiva la
% automática: de otro modo cada sección saldría dos veces y la numeración
% del PDF dejaría de coincidir con la de la app.
\AtBeginSection{}
\AtBeginSubsection{}

\title[Caracterización de un sólido]{Caracterización de un sólido}
\subtitle{Ejemplo didáctico · datos ilustrativos}
\author{Autoría del ejemplo}
\date{\today}

\begin{document}

% --- diapositiva 1 ---
\begin{frame}[plain]
  \transdissolve[duration=0.25]
  \titlepage
\end{frame}
\note{%
  {\scriptsize\textbf{Tiempo previsto: 0:30 min}}\par\medskip
  Ejemplo didáctico. Sustituye contenido y datos por los de tu trabajo. No atribuir estas cifras a un experimento.\par
}

% --- diapositiva 2 ---
\begin{frame}{La pregunta que guía el trabajo}
  \transdissolve[duration=0.25]
  {\large ¿Qué tamaño de cristalita y qué banda prohibida tiene la muestra, y con qué incertidumbre?}
  \medskip

  {\small Caso de uso ilustrativo. Define aquí el sistema, la comparación y el alcance.}
\end{frame}
\note{%
  {\scriptsize\textbf{Tiempo previsto: 1:30 min}}\par\medskip
}

% --- diapositiva 3 ---
\begin{frame}{Una estrategia explícita}
  \transdissolve[duration=0.25]
  \begin{figure}
    \centering
    \definecolor{sa0f}{HTML}{CFDCE6}
    \definecolor{sa1f}{HTML}{9DB7CC}
    \definecolor{sa2f}{HTML}{6B93B2}
    \definecolor{sa3f}{HTML}{396E98}
    \definecolor{sat1843819360}{HTML}{15181C}
    \definecolor{sat1226267613}{HTML}{FFFFFF}
    % Con babel-spanish las palabras largas se parten solas y el texto respira mejor.
    \begin{tikzpicture}[x=1cm,y=1cm,font=\footnotesize]
      \draw[fill=sa0f, draw=none] (0.000,0.000) -- (2.226,0.000) -- (2.518,-1.254) -- (2.226,-2.508) -- (0.000,-2.508) -- cycle;
      \draw[fill=sa1f, draw=none] (2.644,0.000) -- (4.870,0.000) -- (5.162,-1.254) -- (4.870,-2.508) -- (2.644,-2.508) -- (2.936,-1.254) -- cycle;
      \draw[fill=sa2f, draw=none] (5.288,0.000) -- (7.514,0.000) -- (7.806,-1.254) -- (7.514,-2.508) -- (5.288,-2.508) -- (5.580,-1.254) -- cycle;
      \draw[fill=sa3f, draw=none] (7.932,0.000) -- (10.157,0.000) -- (10.450,-1.254) -- (10.157,-2.508) -- (7.932,-2.508) -- (8.224,-1.254) -- cycle;
      \node[text=sat1843819360, align=center, text width=1.75cm, inner sep=1pt] at (1.113,-1.254) {\bfseries Medir un patrón de calibración};
      \node[text=sat1843819360, align=center, text width=1.75cm, inner sep=1pt] at (3.976,-1.254) {\bfseries Ajustar el pico y medir su anchura};
      \node[text=sat1843819360, align=center, text width=1.75cm, inner sep=1pt] at (6.620,-1.254) {\bfseries Aplicar Scherrer con K y λ declarados};
      \node[text=sat1226267613, align=center, text width=1.75cm, inner sep=1pt] at (9.264,-1.254) {\bfseries Transformar la reflectancia y ajustar Tauc};
    \end{tikzpicture}
    \caption{Orden del método: cada paso prepara el siguiente.}
  \end{figure}
\end{frame}
\note{%
  {\scriptsize\textbf{Tiempo previsto: 1:30 min}}\par\medskip
  Explica las decisiones y los controles, no solo los nombres de las técnicas.\par
}

% --- diapositiva 4 ---
\begin{frame}{Mostrar la evidencia}
  \transdissolve[duration=0.25]
  \begin{figure}
    \centering
    \begin{tikzpicture}
      \begin{axis}[
        width=0.90\linewidth,
        height=0.43\linewidth,
        xlabel={2θ (°)},
        ylabel={Intensidad (u. a.)},
        grid=major,
        grid style={line width=.2pt, draw=gray!25},
        tick align=outside,
        tick pos=left,
        axis line style={gray!60},
        label style={font=\small},
        tick label style={font=\footnotesize}
      ]
        \addplot[serieA, line width=1pt, mark=none, smooth] coordinates {
        (20,0) (20.2,0) (20.4,0) (20.6,0) (20.8,0) (21,0) (21.2,0) (21.4,0) (21.6,0) (21.8,0) 
        (22,0) (22.2,0) (22.4,0) (22.6,0) (22.8,0) (23,0) (23.2,0) (23.4,0) (23.6,0) (23.8,0) 
        (24,0) (24.2,0) (24.4,0) (24.6,0) (24.8,0) (25,0) (25.2,0) (25.4,0) (25.6,0) (25.8,0) 
        (26,0) (26.2,0) (26.4,0) (26.6,0) (26.8,0) (27,0) (27.2,0) (27.4,0) (27.6,0) (27.8,0) 
        (28,0) (28.2,0) (28.4,0) (28.6,0) (28.8,0) (29,0) (29.2,0) (29.4,0) (29.6,0) (29.8,0) 
        (30,0) (30.2,0) (30.4,0) (30.6,0) (30.8,0) (31,0) (31.2,0) (31.4,0) (31.6,0) (31.8,0) 
        (32,0) (32.2,0) (32.4,0) (32.6,0) (32.8,0) (33,0) (33.2,0) (33.4,0) (33.6,0) 
        (33.8,0.002) (34,0.013) (34.2,0.063) (34.4,0.21) (34.6,0.5) (34.8,0.841) (35,1) 
        (35.2,0.841) (35.4,0.5) (35.6,0.21) (35.8,0.063) (36,0.013) (36.2,0.002) (36.4,0) 
        (36.6,0) (36.8,0) (37,0) (37.2,0) (37.4,0) (37.6,0) (37.8,0) (38,0) (38.2,0) (38.4,0) 
        (38.6,0) (38.8,0) (39,0) (39.2,0) (39.4,0) (39.6,0) (39.8,0) (40,0) (40.2,0) (40.4,0) 
        (40.6,0) (40.8,0) (41,0) (41.2,0) (41.4,0) (41.6,0) (41.8,0) (42,0) (42.2,0) (42.4,0) 
        (42.6,0) (42.8,0) (43,0) (43.2,0) (43.4,0) (43.6,0) (43.8,0) (44,0) (44.2,0) (44.4,0) 
        (44.6,0) (44.8,0) (45,0) (45.2,0) (45.4,0) (45.6,0) (45.8,0) (46,0) (46.2,0) (46.4,0) 
        (46.6,0) (46.8,0) (47,0) (47.2,0) (47.4,0) (47.6,0) (47.8,0) (48,0) (48.2,0) (48.4,0) 
        (48.6,0) (48.8,0) (49,0) (49.2,0) (49.4,0) (49.6,0) (49.8,0) (50,0) 
        };
      \end{axis}
    \end{tikzpicture}
    \caption{Datos ilustrativos para explicar el método; sustituir por mediciones propias.}
  \end{figure}
  \medskip

  \begin{equation*}
    D=\dfrac{K\lambda}{\beta\cos\theta}
  \end{equation*}
\end{frame}
\note{%
  {\scriptsize\textbf{Tiempo previsto: 1:30 min}}\par\medskip
  Datos o modelos didácticos: reemplazar antes de presentar resultados propios.\par
}

% --- diapositiva 5 ---
\begin{frame}{Interpretación y límites}
  \transdissolve[duration=0.25]
  \begin{itemize}
    \item Pico gaussiano sintético de anchura a media altura (FWHM) de \SI{0.8}{\degree}: con Cu Kα, Scherrer da unos \SI{10}{\nano\meter}. Es un ejemplo de cálculo, no una medida.
    \item Scherrer da un límite inferior si no se descuenta la anchura instrumental; con un patrón de referencia, el valor mejora.
    \item El Eg de Tauc depende del tramo elegido: los tramos vecinos muestran esa sensibilidad.
  \end{itemize}
\end{frame}
\note{%
  {\scriptsize\textbf{Tiempo previsto: 1:30 min}}\par\medskip
}

% --- diapositiva 6 ---
\begin{frame}{Antes de compartir}
  \transdissolve[duration=0.25]
  \begin{itemize}
    \item Anchura instrumental medida con un patrón de referencia
    \item Reflectancia en porcentaje y blanco de referencia declarados
    \item Tramo de Tauc justificado y comparado con los tramos vecinos
  \end{itemize}
  \medskip

  {\small Conserva el proyecto JSON y revisa la exportación final.}
\end{frame}
\note{%
  {\scriptsize\textbf{Tiempo previsto: 1:30 min}}\par\medskip
}

\end{document}